\documentclass[border=10pt]{standalone}
\usepackage[UTF8,fontset=fandol]{ctex}
\usepackage{amsmath,amssymb,tikz}
\usetikzlibrary{arrows.meta,calc,backgrounds}
\definecolor{teal}{HTML}{4F8FA5}
\definecolor{orange}{HTML}{EE995B}
\definecolor{coral}{HTML}{C95B5B}
\definecolor{violet}{HTML}{8A74B5}
\definecolor{muted}{HTML}{656970}
\tikzset{arr/.style={-{Stealth[length=2mm]},draw=black!55,line width=.65pt},
  stage/.style={font=\small\bfseries,text=muted},
  desc/.style={font=\small,text=muted},
  op/.style={font=\Large},
  obj/.style={font=\small},
  dimlabel/.style={font=\scriptsize,text=muted}}
% Matrix face: width = column axis; height = row axis.
% Entries are schematic dense texture, not literal tensor dimensions.
\newcommand{\mat}[8]{%
  \begin{scope}[shift={(#1,#2)}]
  \foreach \r in {1,...,#5}{\foreach \c in {1,...,#6}{
    \pgfmathsetmacro{\tone}{30+mod(floor(abs(sin(13*\r+7*\c))*100),3)*25}
    \fill[#7!\tone,rounded corners=.5pt]
      ({(\c-1)*#3/#6+.015},{-(\r-1)*#4/#5-.015}) rectangle
      ({\c*#3/#6-.015},{-\r*#4/#5+.015});}}
  \draw[black!55,line width=.5pt] (0,0) rectangle (#3,-#4);
  \node[obj,anchor=north] at (#3/2,{-#4-.14}) {#8};
  \end{scope}}
\newcommand{\sh}[3]{\node[dimlabel] at (#1,{#2-.18}) {#3};}

\newcommand{\task}[5]{%
\fill[#4!55,rounded corners=1pt] ({3+(#1)*1.5},#2) rectangle ({3+((#1)+1)*1.5-.04},{#2-.55});
\node[font=\small] at ({3+((#1)+.5)*1.5-.02},{#2-.275}) {#5};}
\begin{document}
\begin{tikzpicture}[x=1cm,y=1cm]
\node[font=\large] at (9.3,1.2) {$T_{\mathrm{serial}}=9\tau,\qquad T_{\mathrm{pipeline}}=5\tau\quad\text{（三块、三阶段的理想例子）}$};
\node[desc] at (9.3,.5) {每块依赖保持为 $D_i\rightarrow C_i\rightarrow R_i$；跨块重叠，不改变单块的先后关系。};
\node[stage,anchor=west] at (0,-.3) {1. 分阶段执行：整批 Dispatch 完成后计算，整批计算完成后回传};
\foreach \row/\lab in {-1.6/Dispatch：接收,-2.55/Compute：专家,-3.5/Return：回传}{
\node[obj,anchor=east] at (2.55,\row-.275) {\lab};
\draw[black!12] (3,\row-.6)--(16.5,\row-.6);}
\foreach \i in {0,1,2}{
\pgfmathtruncatemacro{\j}{\i+1}
\task{\i}{-1.6}{1}{teal}{$D_{\j}$}
\task{\i+3}{-2.55}{1}{orange}{$C_{\j}$}
\task{\i+6}{-3.5}{1}{violet}{$R_{\j}$}}
\draw[arr] (3,-.95)--(16.8,-.95);
\foreach \i in {0,...,9}{\node[dimlabel] at ({3+\i*1.5},-.65) {$\i$};}
\node[dimlabel,anchor=west] at (17.2,-.95) {$t/\tau$};

\node[stage,anchor=west] at (0,-4.65) {2. MegaMoE 的流水线思想：一块到齐就计算，一块算好就回传};
\foreach \row/\lab in {-6.0/Dispatch：接收,-6.95/Compute：专家,-7.9/Return：回传}{
\node[obj,anchor=east] at (2.55,\row-.275) {\lab};
\draw[black!12] (3,\row-.6)--(16.5,\row-.6);}
\foreach \i in {0,1,2}{
\pgfmathtruncatemacro{\j}{\i+1}
\task{\i}{-6.0}{1}{teal}{$D_{\j}$}
\task{\i+1}{-6.95}{1}{orange}{$C_{\j}$}
\task{\i+2}{-7.9}{1}{violet}{$R_{\j}$}}
\draw[arr] (3,-5.35)--(16.8,-5.35);
\foreach \i in {0,...,9}{\node[dimlabel] at ({3+\i*1.5},-5.05) {$\i$};}
\node[dimlabel,anchor=west] at (17.2,-5.35) {$t/\tau$};
\draw[black!45,dashed] (10.5,-5.7)--(10.5,-8.8);
\node[desc,anchor=west,text width=6.2cm] at (11.1,-7.2) {在区间 $[2\tau,3\tau]$：\par 接收第 3 块、计算第 2 块、回传第 1 块。};

\node[anchor=north west,fill=black!4,inner sep=9pt,text width=17.9cm,font=\small,align=left] at (0,-9.4) {
\textbf{维度}\quad 横轴是时间；$i=1,2,3$ 是 token 数据块编号。三个阶段各假设每块耗时 $\tau$。\\[4pt]
\textbf{对象}\quad $D_i$ 接收输入；$C_i$ 包含两次 GEMM 与 SwiGLU；$R_i$ 回传贡献。颜色表示工作类型。\\[4pt]
\textbf{机制}\quad 这是理想教学时间线，不是 profiler 记录。实际存在启动、同步、归并和资源竞争，阶段耗时也不同。};
\end{tikzpicture}
\end{document}
